% GENERATED FILE. Edit canonical lessons, then run npm run generate:books.
% canonical-source-hash: fnv1a64:e421103c3485f64b
% canonical-lessons: FR-C83-soleil, FR-C83-etoile, FR-C83-nuage, FR-C83-pluie, FR-C83-neige

\chapter{Sun, Star, Cloud, Rain, Snow}
\label{ch:fr-sun-star-cloud-rain-snow}
\hlchaptermodality{83}

\begin{chapteropening}
\textbf{By the end of this chapter:} \emph{I can name the sun, a star, a cloud, rain, snow.}

Complete the last lesson of chapter 83: I can name the sun, a star, a cloud, rain, snow.
\end{chapteropening}

\section[le soleil]{le soleil — the sun}

\label{lesson:FR-C83-soleil}



\begin{quote}
Before the new one: say the French for a beach, then the French for an island.
\end{quote}

\subsection*{You'll want to know: le soleil}
\textbf{le soleil} — \textquotedblleft{}the sun\textquotedblright{}.

\begin{cousinweb}
\textbf{soleil} $\leftarrow$ Vulgar Latin \emph{*soliculus}, \textquotedblleft{}a little sun\textquotedblright{}, from Latin \emph{sol}, \textquotedblleft{}sun\textquotedblright{} — inherited. (EtymDB 2.1)
\end{cousinweb}

\begin{grammarlens}[title={What you've built}]
One more word for this chapter.
\end{grammarlens}

\subsection*{Your turn}
\begin{itemize}
\raggedright
  \item \emph{Say it:} \emph{le soleil}
  \item \emph{Say it:} \emph{le soleil}, once more
  \item \emph{Say it:} say \emph{le soleil}
  \item \emph{Recall:} say the French for a beach, then the French for an island, then say \emph{le soleil} again
\end{itemize}

\begin{tcolorbox}[breakable,colback=teal!4,colframe=teal!35!black,title={Before you move on}]
What does le soleil mean? (\textquotedblleft{}The sun\textquotedblright{}.) Say it once more.
\end{tcolorbox}
\section[l'étoile]{l'étoile — a star}

\label{lesson:FR-C83-etoile}



\begin{quote}
Before the new one: say the French for an island, then the French for the sun.
\end{quote}

\subsection*{You'll want to know: l'étoile}
\textbf{l'étoile} — \textquotedblleft{}a star\textquotedblright{}.

\begin{cousinweb}
\textbf{étoile} $\leftarrow$ Old French \emph{estoile} $\leftarrow$ Latin \emph{stella}, \textquotedblleft{}star\textquotedblright{} — inherited. (EtymDB 2.1) The \emph{é-} stands where an \emph{s} once was.
\end{cousinweb}

\begin{grammarlens}[title={What you've built}]
One more word for this chapter.
\end{grammarlens}

\subsection*{Your turn}
\begin{itemize}
\raggedright
  \item \emph{Say it:} \emph{l'étoile}
  \item \emph{Say it:} \emph{l'étoile}, once more
  \item \emph{Say it:} say \emph{l'étoile}
  \item \emph{Recall:} say the French for an island, then the French for the sun, then say \emph{l'étoile} again
\end{itemize}

\begin{tcolorbox}[breakable,colback=teal!4,colframe=teal!35!black,title={Before you move on}]
What does l'étoile mean? (\textquotedblleft{}A star\textquotedblright{}.) Say it once more.
\end{tcolorbox}
\section[le nuage]{le nuage — a cloud}

\label{lesson:FR-C83-nuage}



\begin{quote}
Before the new one: say the French for the sun, then the French for a star.
\end{quote}

\subsection*{You'll want to know: le nuage}
\textbf{le nuage} — \textquotedblleft{}a cloud\textquotedblright{}.

\begin{cousinweb}
\textbf{nuage} — built from Old French \emph{nue}, \textquotedblleft{}a cloud\textquotedblright{}, with \emph{-age}; \emph{nue} $\leftarrow$ Vulgar Latin \emph{*nuba} $\leftarrow$ Latin \emph{nubes}, \textquotedblleft{}cloud\textquotedblright{}. (EtymDB 2.1)
\end{cousinweb}

\begin{grammarlens}[title={What you've built}]
One more word for this chapter.
\end{grammarlens}

\subsection*{Your turn}
\begin{itemize}
\raggedright
  \item \emph{Say it:} \emph{le nuage}
  \item \emph{Say it:} \emph{le nuage}, once more
  \item \emph{Say it:} say \emph{le nuage}
  \item \emph{Recall:} say the French for the sun, then the French for a star, then say \emph{le nuage} again
\end{itemize}

\begin{tcolorbox}[breakable,colback=teal!4,colframe=teal!35!black,title={Before you move on}]
What does le nuage mean? (\textquotedblleft{}A cloud\textquotedblright{}.) Say it once more.
\end{tcolorbox}
\section[la pluie]{la pluie — rain}

\label{lesson:FR-C83-pluie}



\begin{quote}
Before the new one: say the French for a star, then the French for a cloud.
\end{quote}

\subsection*{You'll want to know: la pluie}
\textbf{la pluie} — \textquotedblleft{}rain\textquotedblright{}.

\begin{cousinweb}
\textbf{pluie} $\leftarrow$ Vulgar Latin \emph{*ploia} $\leftarrow$ Latin \emph{pluvia}, \textquotedblleft{}rain\textquotedblright{} (EtymDB 2.1) — inherited: the same root as \emph{il pleut}.
\end{cousinweb}

\begin{grammarlens}[title={What you've built}]
One more word for this chapter.
\end{grammarlens}

\subsection*{Your turn}
\begin{itemize}
\raggedright
  \item \emph{Say it:} \emph{la pluie}
  \item \emph{Say it:} \emph{la pluie}, once more
  \item \emph{Say it:} say \emph{la pluie}
  \item \emph{Recall:} say the French for a star, then the French for a cloud, then say \emph{la pluie} again
\end{itemize}

\begin{tcolorbox}[breakable,colback=teal!4,colframe=teal!35!black,title={Before you move on}]
What does la pluie mean? (\textquotedblleft{}Rain\textquotedblright{}.) Say it once more.
\end{tcolorbox}
\section[la neige]{la neige — snow}

\label{lesson:FR-C83-neige}



\begin{quote}
Before the new one: say the French for a cloud, then the French for rain.
\end{quote}

\subsection*{You'll want to know: la neige}
\textbf{la neige} — \textquotedblleft{}snow\textquotedblright{}.

\begin{cousinweb}
\textbf{neige} — formed from the verb for \textquotedblleft{}to snow\textquotedblright{}, which $\leftarrow$ Vulgar Latin \emph{*nivicare}, from Latin \emph{nix}, \emph{nivis}, \textquotedblleft{}snow\textquotedblright{} (EtymDB 2.1 records the Latin word).
\end{cousinweb}

\begin{grammarlens}[title={What you've built}]
One more word for this chapter.
\end{grammarlens}

\subsection*{Your turn}
\begin{itemize}
\raggedright
  \item \emph{Say it:} \emph{la neige}
  \item \emph{Say it:} \emph{la neige}, once more
  \item \emph{Say it:} say \emph{la neige}
  \item \emph{Recall:} say the French for a cloud, then the French for rain, then say \emph{la neige} again
\end{itemize}

\begin{tcolorbox}[breakable,colback=teal!4,colframe=teal!35!black,title={Before you move on}]
What does la neige mean? (\textquotedblleft{}Snow\textquotedblright{}.) Say it once more.
\end{tcolorbox}
